\section{Exposed paths and a one-particle return identity}
\label{sec:exposure}

The simultaneous test in Proposition~\ref{prop:pin-test} reduces the
probability of many pin-free points to an expected fraction of returns.
Our objective here is to express that first moment through a Markov walk.
We first describe the conditional exchange picture after all pin-free
cycles have been exposed.  A determinant identity then removes the
remaining cycle condition from the first moment.  Finally, we relate the
return probability to the loss of a quadratic form under one period of
the walk.

Throughout this Section, the volume and period are finite.  Let
\(E=D_0\cup E_{\mathrm f}\), where \(E_{\mathrm f}\) is a finite set of
slot-time points.  Let \(T\) be a nonempty set of \(m\) distinct slot points
at one time \(\theta\), disjoint from \(E\), and put \(P=E\cup T\).
Under \(\mathbb P_P\), let
\begin{equation}
 r=\frac1m\sum_{i\in T}
 \mathbb P_P\bigl(\text{the next point of }P\text{ after }i
                         \text{ belongs to }T\bigr).
 \label{eq:exposure-r}
\end{equation}
``After'' always means at strictly positive elapsed time.  In particular,
a return to the same slot point after one or more periods counts.

\subsection{Conditional exchanges on the unexposed slots}

Label every line by its slot at time \(0\), immediately after the
within-site completion.  The one-period endpoint permutation sends these
labels to their slots at the next time \(0\).  Expose the initial labels
on cycles missing \(P\), their trajectories over one period, and their
endpoint mappings.  Write \(\mathcal H\) for the sigma-field generated by
these data and \(\mathbf a\) for an exposure.  At time \(t\), let
\(H_t(\mathbf a)\) be the set of slots not occupied by an exposed label.
We call these slots \emph{holes}.  All pins are holes at their respective
times; in particular every boundary slot is a hole throughout the
period.

The histories are finite-jump paths on a finite state space, together
with finite endpoint data.  Their spaces, and the finite disjoint union
over possible exposed initial sets, are standard Borel spaces.  Thus the
conditional laws below can be taken as regular conditional probabilities.
Equalities of such laws are always asserted almost surely for the
specified outer measure.

For a fixed exposure \(\mathbf a\), define a probability law
\(\mathbb Q_{\mathbf a}\) on the remaining exchange picture as follows.
Independently put rate-\(1/2\) Poisson exchanges on each slot edge during
the time intervals when both endpoints are holes.  At a jump of an
exposed label from a slot to a hole, transport the hole in the opposite
direction, into the vacated slot.  This is a prescribed bijection between
the hole sets just before and just after that jump.  A swap of two
exposed labels leaves the holes unchanged.  At the end of the period,
complete the revealed within-site mappings by independent uniform
bijections of the remaining slots at each site.  The hole sets and the
forced transports are fixed by \(\mathbf a\); the Poisson exchanges and
the completions are the randomness in \(\mathbb Q_{\mathbf a}\).

Let \(G_P\) be the event that every cycle of this remaining picture
meets \(P\), and define
\begin{equation}
 g_P(\mathbf a)=\mathbb Q_{\mathbf a}(G_P).
 \label{eq:exposure-g}
\end{equation}
The set of exposed labels is selected by the condition that their cycles
miss \(P\).  To compute its law, we first follow a deterministic label
set \(A\) under the base law.  Requiring those labels to close among
themselves and avoid \(P\), and every complementary cycle to meet
\(P\), then selects exactly one \(A\) in each complete picture.
The next proposition sums these disjoint possibilities and records both
the conditional law and the exposure marginal used in later averages.

\begin{proposition}[Exposed-path disintegration]
\label{prop:exposure}
For \(P\) as above, the conditional law of the unexposed picture given
\(\mathcal H\) under \(\mathbb P_P\) is
\(\mathbb Q_{\mathbf a}(\,\cdot\mid G_P)\), and
\(g_P(\mathbf a)>0\) for almost every exposure under this pin law.

More precisely, fix a deterministic initial label set
\(A\subseteq V_\Lambda\), and let \(\mu_A\) be the distribution under
\(\mathbb P_{\mathrm b}\) of the histories and endpoint mappings of
these labels.  Let \(F_A\) mean that their endpoint set is \(A\) and
their paths avoid \(P\).  On \(F_A\), let \(c_A(\mathbf a)\) be the
number of cycles in their endpoint permutation.  If \(Z_P\) is the
normalizer defining \(\mathbb P_P\), then for every bounded measurable
function \(\Phi\) of the exposure and remaining picture,
\begin{align}
 &\mathbb E_P\Phi(\mathbf a,\omega_{\mathrm h}) \notag\\
 &\quad=\frac1{Z_P}\sum_{A\subseteq V_\Lambda}
 \int_{F_A}2^{c_A(\mathbf a)}\,
 \mathbb E_{\mathbb Q_{\mathbf a}}
       [\1_{G_P}\Phi(\mathbf a,\omega_{\mathrm h})]\,
                         \mu_A(\dd\mathbf a).
 \label{eq:exposure-disintegration}
\end{align}
Consequently, the exposure marginal on the \(A\)-branch is
\begin{equation}
 \frac1{Z_P}\1_{F_A}(\mathbf a)2^{c_A(\mathbf a)}
                       g_P(\mathbf a)\,\mu_A(\dd\mathbf a).
 \label{eq:exposure-outer-law}
\end{equation}
The same conclusions hold with \(P\) any finite deterministic pin set.
\end{proposition}

\begin{proof}
We first fix \(A\) under the base law and verify the conditional
construction before imposing any cycle condition.  For each slot edge
\(e\), take two independent rate-\(1/2\) Poisson clocks
\(N_e^{(1)},N_e^{(2)}\), all independent across edges.  Accept a mark of
\(N_e^{(1)}\) exactly when at least one endpoint contains an \(A\)-label
immediately before it.  Accept a mark of \(N_e^{(2)}\) exactly when
neither endpoint contains an \(A\)-label.

This construction has the base exchange law.  To see this directly, if
there are \(b\) slot edges, the superposed clock has rate \(b\), and its
next mark has an independent uniform type among the \(2b\) clocks.
For each edge exactly one of its two types is accepted, whatever the
past configuration.  Thus, conditionally on the past, the next
superposed mark is accepted on any specified edge with probability
\(1/(2b)\), and is rejected with probability \(1/2\).  Iterated
conditioning shows that these accepted-edge or rejection outcomes have
the fixed product distribution and are independent of the superposed
waiting times.  Poisson splitting therefore yields independent
rate-\(1/2\) accepted processes on the edges.  If there are no edges the
assertion is immediate.

The motions of the \(A\)-labels before the endpoint permutations use
only the first clock family: second-family marks always exchange two
other labels.  At each site, condition a uniform endpoint permutation on
its images of the \(A\)-labels that are there just before the endpoint.
Each bijection between the remaining source and target slots has the
same number, namely one, of extensions with those prescribed images.
The conditional completion is therefore uniform.  Completions at
different sites remain independent, since the endpoint permutations
were independent and each revealed restriction concerns its own site.

It follows that the full specified histories, including endpoint
images, are determined by the first clock family and these revealed
restrictions, without using the second family or the remaining
completions.  More explicitly, the pre-completion locations of the
$A$-labels depend only on the first clock family.  The endpoint
permutations are independent of both clock families, so conditioning on
these locations and the selected endpoint images leaves the second
clock family independent.  Conditional on each revealed source and
target restriction, the remaining bijection is uniform over its
allowed completions.  Conditional on the histories, the second-family clocks
are still independent Poisson clocks.  Restricting them to the now
fixed intervals with two hole endpoints gives the independent exchanges
in \(\mathbb Q_{\mathbf a}\).  The cross jumps and the remaining
endpoint completions give exactly its other prescribed operations.
Although conditioning on histories may leave unobserved first-family
marks, those marks are rejected whenever they have two hole endpoints
and cannot affect the remaining picture.  This proves the claimed base
conditional kernel on \(F_A\).

On \(F_A\), the labels of \(A\) form a union of whole time cycles
missing \(P\).  Their number is \(c_A(\mathbf a)\), which is determined
by the endpoint mapping of \(A\).  Requiring in addition that every
remaining cycle hit \(P\) makes \(A\) exactly the initial labels of
all the missing cycles.  Therefore for each base picture exactly one
initial set satisfies both requirements.  If \(n_P\) is the number of
cycles missing \(P\), the pointwise identity is
\begin{equation}
 2^{n_P}
   =\sum_{A\subseteq V_\Lambda}
       \1_{F_A}\,2^{c_A}\,\1_{G_P},
 \label{eq:exposure-branch-identity}
\end{equation}
where each summand is defined as zero off \(F_A\).
Apply the base conditional kernel on each branch of this finite sum.
Dividing by \(Z_P=\mathbb E_{\mathrm b}2^{n_P}\) gives
\eqref{eq:exposure-disintegration}.  Integrating the remaining picture
gives \eqref{eq:exposure-outer-law}.  Exposures with
\(g_P(\mathbf a)=0\) consequently have zero outer measure; on its
positive set, division by \(g_P(\mathbf a)\) gives the asserted
conditional law.

For finite pins the same argument applies without change.  For pins
including \(D_0\), avoidance of the boundary is a measurable property
of the finite-jump histories.  Thus the argument applies directly to
that case as well.
\end{proof}

Only the independent law \(\mathbb Q_{\mathbf a}\) will enter the
linear-algebra calculation below.  Conditioning it on \(G_P\) may
destroy the exterior-moment identity for the original layered
permutation.  The factor \(g_P(\mathbf a)\) in
\eqref{eq:exposure-outer-law} also remains present when those conditional
first moments are subsequently averaged under \(\mathbb P_P\).

\subsection{An exterior-power identity and all-hit conditioning}

If \(M\) is a matrix on a finite set \(V\), write
\(\bigwedge^k M\) for its action on the \(k\)-th exterior power; its
entries are the \(k\)-row, \(k\)-column minors of \(M\), with respect
to the chosen ordering of \(V\).  The convention is
\(\bigwedge^0M=1\).  Permutation matrices are column-forward:
\(\Pi e_i=e_{\pi(i)}\).

The all-minors identity below is the determinant-preserving property
studied by Derezi\'nski, Liang and Mahoney
\cite[Definition~4, Example~2, and Lemma~4]{DerezinskiLiangMahoney2020}.
Their rank-one example and closure under independent sums and products
give the finite-product case below.  We include the rank-one argument and the
passage to the exchange law needed here.

\begin{lemma}[Exterior moments before conditioning]
\label{lem:exterior}
Let \(\Pi\) be a random permutation matrix on a finite set \(V\).  Suppose
its law is a limit of laws of finite products of independent Bernoulli
transpositions and deterministic permutations.  Here a Bernoulli
transposition equals a fixed transposition with an arbitrary probability
\(u\in[0,1]\), and equals the identity otherwise.  Then, writing
\(K=\mathbb E\Pi\),
\begin{equation}
 \mathbb E\bigl[\bigwedge\nolimits^k\Pi\bigr]
              =\bigwedge\nolimits^k K
       \qquad (0\le k\le |V|).
 \label{eq:exterior-moments}
\end{equation}

In particular, fix an exposure from Proposition~\ref{prop:exposure}
and finitely many time cuts.  Under \(\mathbb Q_{\mathbf a}\), the
permutation on the disjoint union of the hole layers, obtained by
following each hole from a cut to the next cyclic cut, satisfies
\eqref{eq:exterior-moments}.
\end{lemma}

\begin{proof}
If \(T_{ij}\) transposes two coordinates, then
\[
 T_{ij}-I=-(e_i-e_j)(e_i-e_j)^*.
\]
Every minor of \(I+u(T_{ij}-I)\) is affine in \(u\): a term using
two columns of the rank-one perturbation has determinant zero.
Therefore
\[
 \bigwedge\nolimits^k((1-u)I+uT_{ij})
  =(1-u)\bigwedge\nolimits^kI+u\bigwedge\nolimits^kT_{ij}.
\]
This is \eqref{eq:exterior-moments} for one Bernoulli transposition.
For independent factors \(A,B\) satisfying the identity,
\[
 \mathbb E\bigwedge\nolimits^k(AB)
   =\mathbb E\bigwedge\nolimits^kA\,
      \mathbb E\bigwedge\nolimits^kB
   =\bigwedge\nolimits^k(\mathbb EA\,\mathbb EB)
   =\bigwedge\nolimits^k\mathbb E(AB).
\]
Deterministic permutations satisfy the identity as well, so induction
proves it for finite products.  All the matrix entries and minors are
bounded continuous functions on the finite permutation space, which
proves passage to limits.

We verify the limit representation for the hole picture.  Between
successive exposed jumps and endpoint times, the hole set is fixed and
the available exchanges have constant rates.  Their permutation law has
the Bernoulli-product approximation proved in
Lemma~\ref{lem:pin-exchange-symbol}.
The forced transports are deterministic bijections between equally
sized hole sets.

The same lemma supplies this approximation for every uniform completion.
When its source and target slot sets differ, precede or follow a uniform
permutation of one fixed set by a deterministic bijection.

For completeness, the layered permutation has the same product
structure.  The total number of holes is constant, so identify every
cut layer with one fixed set of that size.  The independent interval
maps act on separate layer blocks; their direct sum, followed by the
fixed cyclic shift of the layers, is the layered permutation.  Each
random operation within a block is a transposition on the full layered
set, and each fixed transport is a permutation there.  Independence
of exchanges on disjoint intervals and of the completions proves the
required representation.
\end{proof}

The return test allows a line started in \(T\) to reach a different
point of \(T\).  If \(\Sigma\) is the column-forward matrix of a
next-pin permutation on a finite pin set containing \(T\), its number
of returns from \(T\) to \(T\) is
\(\sum_{i,j\in T}\Sigma_{ji}\), rather than its trace.
A diagonal-only first-moment identity would therefore not suffice.
The following finite-dimensional statement recovers every entry of
the mean by allowing an arbitrary matrix \(W\).

\begin{proposition}[Conditional next-pin first moment]
\label{prop:first-moment}
Let \(V=J\sqcup Q\) be a finite set with \(J\ne\varnothing\).
Let \(\Pi\) be a random permutation matrix satisfying
\eqref{eq:exterior-moments}, and put \(K=\mathbb E\Pi\).
Let \(G_J\) be the event that every cycle of \(\Pi\) meets \(J\).
On \(G_J\), let \(\Sigma\) be the column-forward permutation matrix
of the next visit to \(J\) at positive time along the cycles.  Then
\begin{equation}
 \mathbb P(G_J)=\det(I-K_{QQ}).
 \label{eq:all-hit-determinant}
\end{equation}
If this probability is positive, \(K_{QQ}\) is transient and
\begin{equation}
 \mathbb E[\Sigma\mid G_J]
   =H:=K_{JJ}+K_{JQ}(I-K_{QQ})^{-1}K_{QJ}.
 \label{eq:conditional-first-moment}
\end{equation}
The matrix \(H\) is the first positive \(J\)-hitting kernel for the
Markov chain with transition matrix \(K\), using independent transitions
on successive steps.  More strongly, for every complex \(J\times J\)
matrix \(W\),
\begin{equation}
 \mathbb E[\det(I-W\Sigma)\mid G_J]=\det(I-WH).
 \label{eq:arbitrary-W}
\end{equation}
An empty \(Q\) is allowed, with determinant of the empty matrix equal
to one and the Schur correction in \eqref{eq:conditional-first-moment}
equal to zero.
\end{proposition}

For time-homogeneous interchange with a single pin, \(G_J\) is the
event that the sampled permutation consists of one cycle.  In that
setting, \eqref{eq:all-hit-determinant} is equivalent to the
single-cycle formula of Alon and Kozma
\cite[Theorem~1]{AlonKozma2013LongCycles}.

\begin{proof}
For a fixed matrix \(D\), the principal-minor expansion and
multiplicativity of exterior powers give
\begin{align*}
 \mathbb E\det(I-D\Pi)
 &=\sum_{k=0}^{|V|}(-1)^k
    \operatorname{Tr}\left(\bigwedge\nolimits^kD\,
                         \mathbb E\bigwedge\nolimits^k\Pi\right)\\
 &=\det(I-DK).
\end{align*}
Take \(D=\operatorname{diag}(W,I_Q)\).  If a permutation cycle lies
entirely in \(Q\), its indicator vector \(v\), extended by zero on
\(J\), satisfies \(D\Pi v=v\), even for arbitrary \(W\).
Thus the determinant on that picture is zero.  On \(G_J\),
\(\Pi_{QQ}\) is nilpotent: a sufficiently long path confined to
\(Q\) would contain a cycle there.  Its inverse resolvent is a finite
sum, and block elimination gives
\begin{align*}
 \det(I-D\Pi)
 &=\det(I-\Pi_{QQ})\,
   \det\left(I-W\left[
      \Pi_{JJ}+\Pi_{JQ}(I-\Pi_{QQ})^{-1}\Pi_{QJ}\right]\right)\\
 &=\det(I-W\Sigma).
\end{align*}
Here \(\det(I-\Pi_{QQ})=1\), and each term of the resolvent follows
the permutation through a specified number of intermediate nonpins.
Consequently the preceding determinant identity reads
\begin{equation}
 \mathbb E[\1_{G_J}\det(I-W\Sigma)]
   =\det\begin{pmatrix}
      I-WK_{JJ}&-WK_{JQ}\\
      -K_{QJ}&I-K_{QQ}
     \end{pmatrix}.
 \label{eq:determinant-before-division}
\end{equation}
At \(W=0\), this proves \eqref{eq:all-hit-determinant} without any
invertibility assumption.  If its value is positive, \(I-K_{QQ}\)
is invertible.  The Schur complement on the right of
\eqref{eq:determinant-before-division}, followed by division by that
positive determinant, gives \eqref{eq:arbitrary-W}.

To extract the first moment, use
\(\det(I-tWA)=1-t\operatorname{Tr}(WA)+O(t^2)\).
The linear terms in \eqref{eq:arbitrary-W} yield
\[
 \operatorname{Tr}\bigl(W\mathbb E[\Sigma\mid G_J]\bigr)
                 =\operatorname{Tr}(WH)
        \qquad\text{for every }W.
\]
Taking the matrix units for \(W\) gives each entry of
\eqref{eq:conditional-first-moment}.

Finally, \(K\) is nonnegative and doubly stochastic, as the mean of
permutation matrices.  Its restriction \(K_{QQ}\) is substochastic.
If a state in \(Q\) could not reach \(J\) by any positive-probability
path, its reachable set would be a nonempty closed subset of \(Q\).
The restriction to that set would have column sums one and eigenvalue
one, forcing \(I-K_{QQ}\) to be singular.  Therefore every state in
\(Q\) has a positive path to \(J\).  Since there are finitely many
states, some integer \(n\) and number \(\varepsilon>0\) make the
probability of leaving \(Q\) within \(n\) steps at least
\(\varepsilon\), uniformly over starting states in \(Q\).  It
follows that
\[
 \|K_{QQ}^{kn}\|_{1\to1}\le(1-\varepsilon)^k,
 \qquad
 (I-K_{QQ})^{-1}=\sum_{j=0}^{\infty}K_{QQ}^j.
\]
The first term of \(H\) records a direct next-step hit of \(J\),
and the term \(K_{JQ}K_{QQ}^jK_{QJ}\) records a first hit after
\(j+1\) intervening steps in \(Q\).  This proves its Markov
interpretation and completes the proof.
\end{proof}

We have now identified the conditional first moment at a finite set of
cuts.  In the hole construction, each application of \(K\) in this
Markov interpretation resamples the relevant interval transition.
It does not follow a single sampled permutation repeatedly.  This
distinction permits the return calculation to use one-particle heat
flow in the fixed exposed environment.  Figure~\ref{fig:exposure-first-moment}
summarizes the finite-cut reduction and the two probability laws it
distinguishes.

\begin{figure}[tb]
\centering
\begin{tikzpicture}[
  box/.style={draw=blue!45!black,rounded corners=2pt,fill=blue!3,
              align=center,text width=5.2cm,minimum height=1.3cm,inner sep=6pt},
  >={Stealth[length=2mm]},font=\small]
  \node[box] (independent) at (0,0)
    {Unconditioned hole picture $\mathbb Q_{\mathbf a}$\\
     $K=\mathbb E_{\mathbb Q_{\mathbf a}}\Pi$};
  \node[box] (conditioned) at (7,0)
    {Conditioned hole picture\\
     $\mathbb Q_{\mathbf a}(\,\cdot\mid G_J)$};
  \node[box] (walk) at (0,-2.6)
    {Markov walk with transition $K$\\
     First positive $J$-hitting kernel};
  \node[box] (moment) at (7,-2.6)
    {Conditional next-pin mean\\
     $\mathbb E[\Sigma\mid G_J]$};
  \draw[->,thick] (independent) -- node[above,align=center]
    {condition\\on $G_J$} (conditioned);
  \draw[->,thick] (independent) -- node[left,align=right]
    {retain the\\one-step mean} (walk);
  \draw[->,thick] (conditioned) -- node[right,align=left]
    {retain the\\next-pin mean} (moment);
  \draw[<->,thick] (walk) -- node[below,align=center]
    {equal\\kernels} (moment);
\end{tikzpicture}
\caption{The first-moment reduction at a fixed exposure and a finite set $J$
of pins on selected time cuts, with $\mathbb Q_{\mathbf a}(G_J)>0$.  The exterior-moment identity is used for the unconditioned
law in the upper-left box; Proposition~\ref{prop:first-moment} proves
that the two routes to the lower-right box agree.  The lower-left walk
resamples its transitions, whereas the upper-right law follows a single
sampled permutation.  The equality concerns the next-pin mean and does
not identify these two path laws.  When the exposure is averaged, its
actual marginal still contains the factor $g_P(\mathbf a)$ in
\eqref{eq:exposure-outer-law}.}
\label{fig:exposure-first-moment}
\end{figure}

\FloatBarrier

\subsection{Fresh turns and continuous boundary pins}

Given an exposure \(\mathbf a\) with \(g_P(\mathbf a)>0\), define
the \emph{fresh-turn walk} on holes as follows.  Repeat periodically the
hole sets, forced transports, and revealed endpoint restrictions.
Between those prescribed changes, the walk jumps across every available
slot edge at rate \(1/2\).  At a forced transport it follows that
transport; at an endpoint completion it uses an independent uniform
completion at the current site.  Use fresh Poisson exchanges and
completion choices in each turn.  These prescriptions define a
time-inhomogeneous Markov walk.  On every bounded time interval it has
finitely many jumps almost surely.  For a starting point \(i\in T\),
write \(\tau_T^+\) for its first strictly positive visit to \(T\),
and \(\tau_E\) for its first visit to \(E\), allowing infinite
hitting times.

\begin{lemma}[Return identity with boundary pins]
\label{lem:exposure-dense-return}
For \(\mathbb P_P\)-almost every exposure and every \(i\in T\),
\begin{equation}
 \mathbb P_P\bigl(\text{next pin after }i\text{ lies in }T
                              \mid\mathcal H\bigr)
   =\mathbb P_{\mathbf a}^{\mathrm{walk}}
                         (\tau_T^+<\tau_E).
 \label{eq:fresh-return-identity}
\end{equation}
In particular the first moment in \eqref{eq:exposure-r} is the
\(\mathbb P_P\)-average of the corresponding fresh-turn return
probabilities in the exposed environment.
\end{lemma}

\begin{proof}
Fix an exposure for which Proposition~\ref{prop:exposure} holds and
\(g_P(\mathbf a)>0\).  Choose increasing finite deterministic time
sets \(\mathcal T_n\) with dense union, and let \(D_{0,n}\) contain
all boundary slots at the times in \(\mathcal T_n\).  Put
\[
 E_n=E_{\mathrm f}\cup D_{0,n},\qquad P_n=E_n\cup T.
\]
The cuts used below include all times occurring in \(P_n\), as well
as time \(0\) if necessary.  An extra unpinned cut merely subdivides
an interval and does not change which cycles hit \(P_n\).

Under \(\mathbb Q_{\mathbf a}\), let \(G_n\) be the event that
all hole cycles meet \(P_n\), and set
\(g_n=\mathbb Q_{\mathbf a}(G_n)\).  The events \(G_n\) increase.
Almost surely the one-period hole picture has only finitely many
random exchanges, in addition to the finitely many prescribed changes.
Its spatial visits have positive residence time: independent Poisson
jumps do not coincide with any prescribed change, and within-site
completions do not change the site.  The exposed histories avoid
\(D_0\), so their cross jumps cannot create an instantaneous
boundary excursion either.  Every boundary visit by a hole cycle is
therefore detected by some cut in the dense union.  There are finitely
many cycles, whence
\begin{equation}
 \1_{G_n}\uparrow\1_{G_P}\quad
       \mathbb Q_{\mathbf a}\text{-almost surely},\qquad
 g_n\uparrow g_P(\mathbf a)>0.
 \label{eq:dense-all-hit-limit}
\end{equation}
In particular, \(g_n>0\) for all sufficiently large \(n\).
No invertibility is required for the earlier indices.

For such an \(n\), Lemma~\ref{lem:exterior} applies to the
unconditioned layered hole permutation.  With \(J=P_n\),
Proposition~\ref{prop:first-moment} identifies each next-pin probability
under \(\mathbb Q_{\mathbf a}(\,\cdot\mid G_n)\) with the
corresponding first-hit probability for its matrix \(K\).  The latter
is precisely the walk just defined, observed at the finitely many cuts:
interval choices are independent, including on later passes through
the same interval.  Thus, if \(A_n\) is the event in a fixed hole
picture that the next \(P_n\)-pin after \(i\) belongs to \(T\),
\begin{equation}
 \frac{\mathbb Q_{\mathbf a}(G_n\cap A_n)}{g_n}
      =\mathbb P_{\mathbf a}^{\mathrm{walk}}
                         (\tau_T^+<\tau_{E_n}).
 \label{eq:finite-cut-return-ratio}
\end{equation}

On the picture side, a cycle through \(i\in T\) returns to \(T\)
after finitely many periods, regardless of which other pins it meets.
Detection of boundary visits before that return shows that
\(\1_{A_n}\) decreases to the event \(A\) that the next full
\(P\)-pin is in \(T\).  Together with
\eqref{eq:dense-all-hit-limit} and bounded convergence, this gives
\[
 \mathbb Q_{\mathbf a}(G_n\cap A_n)
       \longrightarrow\mathbb Q_{\mathbf a}(G_P\cap A).
\]
Dividing by the positive limiting denominator yields the return
probability under \(\mathbb Q_{\mathbf a}(\,\cdot\mid G_P)\).

For the walk side, use a single uncensored fresh-turn trajectory for
all the tests.  Its first positive \(T\)-hitting time \(\tau_T^+\)
is then fixed, independently of \(n\).  The events on the right of
\eqref{eq:finite-cut-return-ratio} decrease.  A path in their
intersection has \(\tau_T^+<\infty\), since this is already part
of each event.  If that path had visited the boundary before
\(\tau_T^+\), the positive residence time of the visit would make
one of the dense cuts detect it before \(\tau_T^+\).  Thus their
intersection is exactly \(\{\tau_T^+<\tau_E\}\), up to a null
set.  Continuity from above gives convergence of the probabilities.
Finally Proposition~\ref{prop:exposure} identifies the limiting
conditional picture law with the conditional law under
\(\mathbb P_P\).  This proves \eqref{eq:fresh-return-identity}.
\end{proof}

\subsection{Minimal hitting probabilities and return energy}

We now ignore all pins of \(E\) except \(D_0\) in the stopping
rule for this walk.  The exposed environment is kept fixed.  This can
only increase its probability of returning to \(T\).  Let
\(U=U(\mathbf a)\) be its one-period column-forward transition
matrix from the cut at \(\theta\) to that same cut one period later,
killed upon every visit to the boundary, including at its initial
point.  Boundary coordinates can be retained with value zero.
With the post-jump convention at deterministic cuts, a period includes
all changes after its starting cut and through its ending cut; in
particular the within-site completion is included exactly once.
All subsequent periods use fresh randomness.

We record the basic contraction property for later heat-flow estimates.
For any subinterval, a realization of all exchanges, forced transports,
and completions gives a bijection between its initial and final hole
sets.  Keeping only trajectories that avoid the boundary gives a
partial permutation: its matrix has at most one entry equal to one
in each row and in each column.  The killed transition matrix \(C\)
is the mean of these matrices.  Hence
\begin{equation}
 C_{ij}\ge0,\qquad
 \sum_i C_{ij}\le1,\qquad \sum_j C_{ij}\le1,
 \qquad \|Cv\|_2\le\|v\|_2.
 \label{eq:hole-contractions}
\end{equation}
Indeed Cauchy--Schwarz in each row gives
\(\sum_i|(Cv)_i|^2\le\sum_{i,j}C_{ij}|v_j|^2
\le\sum_j|v_j|^2\).  The same inequalities apply to adjoints and
to additional killing.  Row sums bound propagated values by one;
column sums bound the total mass from a unit starting point.  These
are statements at a fixed exposure, without any assertion of symmetry
of \(U\) or invariance of the exposure law under time reversal.

\begin{lemma}[Minimal hitting function and return energy]
\label{lem:return-energy}
For almost every exposure under \(\mathbb P_P\), let \(U\) be the
boundary-killed one-period matrix just defined.  There is a smallest
nonnegative function \(f\) on its cut holes satisfying
\begin{equation}
 f=1\text{ on }T,\qquad U^*f=f\text{ off }T.
 \label{eq:minimal-hitting-harmonicity}
\end{equation}
It is given off \(T\) by the probability of ever reaching \(T\)
on a future turn before boundary killing.  In particular
\(0\le f\le1\), and \(f\) is zero on the boundary.  With counting
inner products,
\begin{equation}
 \langle f,(I-U)f\rangle
        =\sum_{i\in T}\bigl(1-(U^*f)_i\bigr),
 \qquad
 m(1-r)\ge
       \mathbb E_P\langle f,(I-U)f\rangle.
 \label{eq:return-energy}
\end{equation}
Both \(U\) and \(f\) in this expectation are functions of the same
exposure, with the outer law in \eqref{eq:exposure-outer-law}.
\end{lemma}

\begin{proof}
Set \(f_0=\1_T\) and recursively define
\begin{equation}
 f_{n+1}=\1_T+\1_{H_\theta\setminus T}\,U^*f_n,
 \label{eq:finite-horizon-hitting}
\end{equation}
where the multiplication by the indicator is pointwise.
Nonnegativity and the substochasticity of \(U^*\) show inductively
that \(0\le f_n\le f_{n+1}\le1\).  The Markov property identifies
\(f_n\), off \(T\), with the probability of reaching \(T\)
within the next \(n\) periods before killing.  Thus its pointwise
limit \(f\) has the asserted probability interpretation.  There are
finitely many coordinates, so passage to the limit in
\eqref{eq:finite-horizon-hitting} gives
\eqref{eq:minimal-hitting-harmonicity}.  A boundary starting point is
killed immediately, and consequently every \(f_n\) and \(f\) is
zero there.

If \(h\) is any nonnegative solution of
\eqref{eq:minimal-hitting-harmonicity}, then \(h\ge f_0\) and
induction in \eqref{eq:finite-horizon-hitting} gives \(h\ge f_n\)
for every \(n\); hence \(h\ge f\).  This proves minimality.
In particular \(f\) is zero on any closed class that cannot reach
\(T\).  No assumption of certain absorption by the boundary is
needed.

Since \(T\) occurs at a single time modulo the period, a strictly
positive return from that cut to \(T\) occurs after an integer number
of periods.  Conditioning on the first period therefore shows that
the relaxed return probability from \(i\in T\) is \((U^*f)_i\).
By Lemma~\ref{lem:exposure-dense-return}, the conditional probability
\(q_i(\mathbf a)\) contributing to \eqref{eq:exposure-r} is the
return probability with all of \(E\) retained.  Hence
\[
 q_i(\mathbf a)\le (U^*f)_i\qquad(i\in T).
\]
The vectors are real, so transposition in the scalar quadratic form
and harmonicity off \(T\) give
\begin{align*}
 \langle f,(I-U)f\rangle
   &=\langle (I-U^*)f,f\rangle\\
   &=\sum_{i\in T}\bigl(1-(U^*f)_i\bigr)
    \le\sum_{i\in T}(1-q_i(\mathbf a)).
\end{align*}
Averaging under the actual exposure marginal gives
\eqref{eq:return-energy}.  Measurability causes no additional
limiting issue: finite-time transition probabilities are measurable
functions of the finite-jump exposure, and \(f\) is the increasing
limit of the finite-horizon functions in
\eqref{eq:finite-horizon-hitting}.
\end{proof}

For the estimates that follow, write \(\eta=\1_T\) and \(g=f-\eta\).
Then \(0\le g\le1\), \(\langle\eta,g\rangle=0\), and both
functions have zero boundary values.  Lemma~\ref{lem:return-energy}
shows that the desired scale of quadratic loss is \(m\): for any
\(\varepsilon\ge0\),
\[
 \mathbb E_P\langle f,(I-U)f\rangle\ge m(1-\varepsilon)
 \quad\Longrightarrow\quad r\le\varepsilon.
\]
The simultaneous pin test then bounds the old-pin avoidance probability
by \((2\varepsilon)^m\).  The next sections obtain this lower bound
with a small error from the spatial distribution of holes and the
corresponding heat flow.
